\subsection{Independent continuous cup vanishing for ordinary pairs}
\label{sec:ghost-ordinary-cup}

\begin{proposition}[The ordinary continuous cup vanishes]
\label{prop:ghost-ordinary-cup}
Every ordinary completed pair \((A,B)\in\mathcal G_{\rm ord}\)
satisfies
\[
           A\cup B=0\quad\hbox{in}\quad
 H^2_{\rm cts}(G_{\mathbb Q},\widehat{\mathbb Z}(2)).
\]
This proves the continuous cup assertion used in SECTIONS A.27
without assuming its omitted strict localization comparison.
\end{proposition}
\begin{proof}
Fix \(\ell\), put \(T=\mathbb Z_\ell(2)\),
\(T_n=T/\ell^nT\), and choose literal integral Kummer
cocycles \(a,b\) for the two coordinates. We use additive
notation and write \(\chi\) for the cyclotomic character.

First, the class of \(a\cup b\) is zero at every finite
coefficient exponent. Here is the exact global detection
argument from SECTIONS A.25. The tame-symbol sequence is
\[
 0\longrightarrow K_2(\mathbb Z)\longrightarrow K_2(\mathbb Q)
 \longrightarrow\bigoplus_q\mathbb F_q^\times\longrightarrow0,
 \qquad K_2(\mathbb Z)=\langle\{-1,-1\}\rangle=\mathbb Z/2.
\]
The Galois symbol identifies \(K_2(F)/m\) with
\(H^2(G_F,\mu_m^{\otimes2})\). For an element locally
zero modulo \(m\), every tame coordinate is divisible by
\(m\). Lift a finitely supported vector of these divisors
to \(K_2(\mathbb Q)\); subtraction leaves an element of
\(K_2(\mathbb Z)\). For odd \(m\) it is zero modulo
\(m\); for even \(m\) its nonzero element is detected by
the quadratic Hilbert symbol at two. Indeed
\(X^2+Y^2=-1\) has no \(\mathbb Q_2\)-solution: clearing
denominators gives a primitive sum of three squares zero
modulo four, which is impossible. Thus finite-prime
localization in this \(H^2\) is injective. The finite local
cup of \((t_q,1-t_q)\) is zero by the Steinberg relation,
so the global finite cup is zero. The exact algebraic inputs
are \cite[III.5.2.2, III.6.5--6.5.1, III.6.10.3--6.10.4]{GhostWeibelKbook}.

Put \(S=\{\ell,\infty\}\). The weight-two residue map
\[
 H^1(G_{\mathbb Q},T_n)\longrightarrow
       \bigoplus_{q\ne\ell}H^0(\mathbb F_q,T_n(-1))
       \tag{\dagger}
\]
is surjective at every \(n\). More explicitly, for
\(X=\operatorname{Spec}\mathbb Z[1/\ell]\), remove a finite
set of closed points and take the cohomology sequence with
supports. At a removed point \(q\), the formula
\(R^2i^!\mu_{\ell^n}=\mathbb Z/\ell^n\), with the other
terms zero, gives
\(H^2_q(X,T_n)=H^0(\mathbb F_q,T_n(-1))\).
This is \cite[II.1.7(b)]{GhostMilneADT}, tensored with the
locally free sheaf \(\mu_{\ell^n}\). Taking the filtered
union over finite sets, and applying the group-to-scheme
comparison \cite[II.2.9]{GhostMilneADT}, gives the exact
sequence
\[
 H^1(G_{\mathbb Q},T_n)\longrightarrow
 \bigoplus_{q\ne\ell}H^0(\mathbb F_q,T_n(-1))
 \longrightarrow H^2(G_{\mathbb Q,S},T_n)
 \longrightarrow H^2(G_{\mathbb Q},T_n).
\]
The first arrow is precisely restriction to inertia followed
by the tame-character identification. Thus its cokernel is
the kernel of the last arrow.
For odd \(\ell\), the first group is zero: its integral
version is \(K_2(\mathbb Z)\{\ell\}=0\), and integral
degree three is zero as proved there. For \(\ell=2\),
\cite[VI.9.11, case \(n=2\)]{GhostWeibelKbook} gives
\(H^2(G_{\mathbb Q,S},\mathbb Z_2(2))\simeq\mathbb Z/2\).
Its real restriction is an isomorphism. To verify detection
directly, the cocycle
\(\epsilon(g)=(\chi(g)-1)/2\) represents \(\kappa(-1)\),
and \(\epsilon\cup\epsilon\) at two complex conjugations
has the odd value \(-1\), generating real \(H^2\).
Integral degree three is zero: finite cohomology in degree
three is its real counterpart by
\cite[I.4.10(c)]{GhostMilneADT}, and reduction sends its
order-two generator \(2^n\) at level \(n+1\) to zero
at level \(n\). The finite degree-two groups have zero
derived limit. These are ordinary positive-degree real
cohomology groups: for the trivial real action on \(T_n\),
\(H^3(G_{\mathbb R},T_n)=T_n[2]\), whereas
\(H^2(G_{\mathbb R},T_n)=T_n/2T_n\).
They coincide with Tate cohomology in the positive degrees
used here. No assertion uses ordinary degree-zero cohomology
in place of Tate degree zero. Thus the coefficient exact sequence makes
\(H^2(G_{\mathbb Q,S},T_n)\to H^2(G_{\mathbb R},T_n)\)
an isomorphism for every \(n\), so its map to absolute
cohomology is injective. This proves \((\dagger)\) also
at two.

We next choose integral local primitives at every
\(q\ne\ell\) with controlled inertia values. Let
\(a_q=v_q(A)\), \(b_q=v_q(B)\), ordinary integers.
The local Kummer character of a principal \(q\)-unit is
zero at coefficient prime \(\ell\), since all its
\(\ell^n\)-th roots lie in the principal unit group.
Consequently:
\begin{itemize}
 \item If \(a_q>0\), then \(b=\delta u\) locally for
 some \(u\in\mathbb Z_\ell(1)\). Take
 \(h_q(g)=-a(g)\chi(g)u\).
 \item If \(b_q>0\), then \(a=\delta u\) locally.
 Take \(h_q(g)=u b(g)\).
 \item If \(a_q=b_q<0\), the equality
 \((1-t_q)/t_q=-(1-1/t_q)\) gives
 \(b=a+\epsilon+\delta u\), where
 \(\epsilon(g)=(\chi(g)-1)/2\) represents
 \(\kappa(-1)\). Take
 \[
 h_q(g)=-\frac{a(g)^2+\chi(g)a(g)}2
                   -a(g)\chi(g)u.
 \]
 This expression is integral even at two: \(\chi(g)\)
 is odd, so \(a(g)\) and \(a(g)+\chi(g)\) have opposite
 parity. Direct expansion gives
 \(\delta[-(a^2+\chi a)/2]=a\cup(a+\epsilon)\), and
 \(\delta[-a\chi u]=a\cup\delta u\).
 \item If \(a_q=b_q=0\), the literal cocycles are
 unramified. Their cup on the unramified procyclic
 quotient has an integral primitive; choose that
 primitive inflated to the local group.
\end{itemize}
In every case \(\delta h_q=a\cup b\). On inertia
\(\chi=1\), so all values of \(h_q|_{I_q}\) belong
to \(a_q\mathbb Z_\ell\) or \(b_q\mathbb Z_\ell\)
in the first two cases, to
\(\ell^{-v_\ell(2)}a_q\mathbb Z_\ell\) in the third,
and are zero in the fourth. Both global valuation
vectors tend to zero \(\ell\)-adically. Therefore for
every \(n\), these prescribed inertia values are zero
modulo \(\ell^n\) at all but finitely many \(q\).

Choose a global finite primitive \(f_n\) of the cup
modulo \(\ell^n\), whose existence was proved above.
For every \(q\ne\ell\), the difference
\(f_n|_{G_{\mathbb Q_q}}-h_q\bmod\ell^n\) is a
local one-cocycle. Its tame restriction gives an
element of \(H^0(\mathbb F_q,T_n(-1))\). These elements
have finite support: a continuous finite-valued global
cochain factors through a finite Galois quotient and
vanishes on inertia outside finitely many primes after
normalization, and the prescribed restrictions have
the same finite-support property. Surjectivity in
\((\dagger)\) gives a global one-cocycle with exactly
these residues. Inertia acts trivially on \(T_n\) at
\(q\ne\ell\), so its coboundaries are zero there; equality
of the residue classes is equality of the actual inertia
homomorphisms. Subtract this global cocycle from \(f_n\).

Thus the set \(P_n\) of global finite primitives
whose inertia restrictions equal those of the fixed
\(h_q\) at every \(q\ne\ell\) is nonempty. Differences
of its elements are precisely the one-cocycles
vanishing on all outside-\(S\) inertia, hence factoring
through \(G_{\mathbb Q,S}\). Therefore \(P_n\) is a
torsor under the finite group \(Z^1(G_{\mathbb Q,S},T_n)\).
The coefficient maps preserve these exact restrictions,
so the nonempty finite \(P_n\) form an inverse system.
Every finite set of compatibility equations is solved
by projection from its highest level; compactness gives
a compatible family. Its limit is a continuous
integral primitive of \(a\cup b\).

Finally take the product of these primitives over all
\(\ell\). A map into the product is continuous exactly
when its components are continuous, and the cochain
differential is coordinatewise. This gives the required
\(\widehat{\mathbb Z}(2)\)-valued primitive.
\end{proof}

\begin{corollary}[Every ordinary pair has an actual global class-two lift]
\label{cor:ghost-class2-lift}
Put \(U=\mathbb P^1_{\mathbb Q}\setminus\{0,1,\infty\}\),
base it at the tangent vector \(\overrightarrow{01}\), and
write \(\Pi=\pi_1(U_{\overline{\mathbb Q}},\overrightarrow{01})\).
Every ordinary completed pair is the abelianization of a
continuous section of
\[
 \pi_1(U,\overrightarrow{01})/\Gamma_3\Pi
       \longrightarrow G_{\mathbb Q}.
\]
Consequently no ordinary pair is excluded by unrestricted
global lifting at nilpotency class one or two.
\end{corollary}
\begin{proof}
Let \(s_0\) be the tangential section, and choose the usual
small positive loops \(x,y\) around zero and one, with the
real interval used to transport the second loop. In these
literal generators the arithmetic action is
\[
 g(x)=x^{\chi(g)},\qquad
 g(y)=\varphi_g^{-1}y^{\chi(g)}\varphi_g,
             \qquad\varphi_g\in\Gamma_2\Pi;
\]
see \cite[(12.11)--(12.15)]{GhostWickelgren2012} for these
base points, generators, and action. In \(N_2=\Pi/\Gamma_3\Pi\)
the commutator subgroup is central, so this action sends
\(y\) to \(y^{\chi(g)}\). With the convention
\(C=[x,y]=xyx^{-1}y^{-1}\), every element is uniquely
\(y^b x^a C^h\), and multiplication and action are
\[
 (a,b,h)(a',b',h')=(a+a',b+b',h+h'+ab'),\qquad
 g(a,b,h)=(\chi a,\chi b,\chi^2h).
\]
These identities first follow by discrete collection and
then hold in \(\widehat{\mathbb Z}^3\) by continuity,
as in Lemma~\ref{lem:ghost-hall-model}. The abelianization
sends a point \(t\) to \((\kappa(t),\kappa(1-t))\),
so the specified completed pair has literal Kummer
representatives \(a,b\) in exactly these coordinates.
By Proposition~\ref{prop:ghost-ordinary-cup}, choose a
continuous \(f:G_{\mathbb Q}\to\widehat{\mathbb Z}(2)\)
with \(\delta f=a\cup b\), where
\((a\cup b)(g,h)=a(g)\chi(g)b(h)\). Then
\[
 s(g)=y^{b(g)}x^{a(g)}C^{-f(g)}s_0(g)
\]
is a continuous section: the displayed multiplication law
reduces its homomorphism identity exactly to
\(\delta(-f)=-a\cup b\). This constructs the section
in the actual arithmetic quotient, including all coefficient
primes. Its abelianization is the original pair.
\end{proof}
